\section{Background}
\label{sec:background}

\subsection{Centred coordinates and the height function}
\label{sec:bg_coords}

Let $S=(s_{ij})$ be an arrangement of order $n$, with cells indexed by $i,j\in\{0,\dots,n-1\}$ ($i$ the row from the top, $j$ the column from the left), and write $N=n^2$.

\begin{definition}[Centred coordinates and height function]
\label{def:coords}
To the cell $(i,j)$ assign
\begin{equation}
x_{ij}=j-\frac{n-1}{2},\qquad y_{ij}=\frac{n-1}{2}-i,\qquad z_{ij}=s_{ij}-\frac{N+1}{2}.
\label{eq:coords}
\end{equation}
The function $z=(z_{ij})$ on the grid is the \emph{height function} of $S$.
\end{definition}

The horizontal coordinates centre the grid at the origin; the vertical coordinate centres the entries so that the mean entry $(N+1)/2$ sits at height zero.
All three have mean zero over the $N$ cells (\Cref{lem:variances}).
The $N$ points $(x_{ij},y_{ij},z_{ij})\in\R^3$ have a convex hull, the object of the companion paper~A, \emph{Magic Gems: A Polyhedral Framework for Magic Squares} (Mathewson, arXiv:2512.09170) \cite{mathewson_gems}, which studies the convex hull of these points, the symmetry classification of the $880$ order-four hulls, and the hulls of Siamese squares.
Nothing below uses the hull.

\subsection{Covariance and Loly's planar moment}
\label{sec:bg_cov}

For functions $\varphi,f$ on the cells, $\Cov(\varphi,f)$ denotes the empirical covariance over the $N$ cells, $\Cov(\varphi,f)=\frac1N\sum_c(\varphi_c-\bar\varphi)(f_c-\bar f)$, and $X,Y,Z$ denote the functions \eqref{eq:coords}.
Loly's planar moment of inertia \cite{loly2004invariance} places the \emph{value} $s_{ij}$ as a mass at the planar point $(x_{ij},y_{ij})$ and computes $\sum s_{ij}(x_{ij}^{2}+y_{ij}^{2})$.
It is an affine function of the single degree-two cross-moment $\Cov(Z,X^2+Y^2)$ (\Cref{cor:loly}).

\subsection{Two tools from permutation statistics}
\label{sec:bg_perm}

Let $\pi$ be a uniformly random permutation of $\{1,\dots,N\}$ and $a,b\in\R^N$ fixed vectors.
The quantity $T_a=\sum_i a_i b_{\pi(i)}$ is a \emph{linear permutation statistic}.
Its mean is $N\bar a\bar b$ and, for two such statistics with the same $b$ (Wald--Wolfowitz \cite{waldwolfowitz1944}, polarised),
\begin{equation}
\Cov(T_a,T_{a'})=\frac{1}{N-1}\sum_i(a_i-\bar a)(a'_i-\bar a')\sum_j(b_j-\bar b)^2 .
\label{eq:ww_identity}
\end{equation}
Hoeffding's combinatorial central limit theorem \cite[Thm.~3]{hoeffding1951} gives asymptotic normality of $\sum_i c_N(i,\pi(i))$ when, for every integer $r\ge3$, $N^{-1}\sum_{i,j}|d_N(i,j)|^r$ is $O\bigl((N^{-1}\sum_{i,j}d_N(i,j)^2)^{r/2}\bigr)$, $d_N$ being the doubly centred array; for product arrays $c_N(i,j)=a_ib_j$ the condition reduces to a bound on the standardised absolute moments of $a$ and $b$.
In \Cref{sec:random_energy} the energies of a random arrangement are quadratic forms in such statistics; \eqref{eq:ww_identity} gives their means exactly, and the distribution of a quadratic form in Gaussian variables can be computed by numerical inversion of its characteristic function, as in Imhof's method \cite{imhof1961}.

\subsection{Enumeration and related polytopes}
\label{sec:bg_polytopes}

There are $8$ magic squares of order $3$, all images of the Lo Shu under the dihedral group $D_4$ of the square; $7{,}040$ of order $4$, forming $880$ $D_4$-orbits first listed by Fr\'enicle de Bessy and verified many times since \cite{beck2010enumeration}; and $275{,}305{,}224$ of order $5$ \cite{gardner1976}.
For $n\ge6$ only statistical estimates are known \cite{trump2012enumeration}.
Classical constructions show that magic squares exist for all $n\ge3$ \cite{andrews1960magic}.

The polytope literature concerns a different object from the one used here.
An $n\times n$ nonnegative integer matrix with all row and column sums equal to $t$ is a lattice point of the $t$-th dilate of the Birkhoff polytope, and their number is a polynomial in $t$ \cite{stanley1973magic,ehrhart1973magic,ziegler1995lectures}; adding the diagonal conditions and distinct entries leads to the inside-out polytopes of Beck and Zaslavsky \cite{beck2006insideout,beck2010enumeration}; the cones of magic squares and cubes have Hilbert-basis descriptions \cite{ahmed2003polyhedral}; and contingency tables with fixed margins are the same lattice points studied as a sample space \cite{diaconis1995rectangular}.
Spectral and matrix properties of a magic square \cite{loly2009spectra,nordgren2012properties,nordgren2014spectra}, the singular-value signatures of Cameron, Rogers and Loly \cite{cameron2013signatura}, and Loly's physical invariants \cite{loly2004invariance,rogers2004inertia,rogers2005multipole,loly2007franklin} are closer to our viewpoint: they attach a quantity to one square and ask what it reveals.
